\chapter{Active Circuit Building Blocks}
\section{Why Active Circuits Appear}
The Extra exam includes amplifiers, op-amps, feedback, oscillators, mixers, and
regulated power supplies. A full electronics text is not needed here; the goal
is to connect the circuit formulas to system behavior.

\section{Amplifier Gain}
Voltage gain is
\[
A_v=\frac{V_o}{V_i}.
\]
In decibels,
\[
A_{v,\dB}=20\log_{10}|A_v|
\]
when the input and output voltage ratio is the relevant quantity. Power gain is
\[
G_{\dB}=10\log_{10}\frac{P_o}{P_i}.
\]

\section{Op-Amp Ideals}
An ideal op-amp has very high input impedance, very low output impedance, and
very large open-loop gain. Negative feedback makes practical closed-loop gains
predictable.

For an inverting amplifier,
\[
\frac{V_o}{V_i}=-\frac{R_f}{R_{in}}.
\]
For a non-inverting amplifier,
\[
\frac{V_o}{V_i}=1+\frac{R_f}{R_g}.
\]

\section{Feedback}
Negative feedback trades gain for bandwidth, linearity, reduced sensitivity to
device variation, and often lower distortion. It can also create stability
questions if phase lag becomes large near unity loop gain.

\begin{controlsbox}
Negative feedback is a tradeoff involving gain, bandwidth, sensitivity,
distortion, and stability. That statement is exactly the same in RF amplifiers
and process-control loops.
\end{controlsbox}

\section{Oscillators}
Oscillation requires sustained positive-feedback behavior at a frequency where
the loop phase and loop gain conditions support a non-decaying sinusoid. In
practice, nonlinear amplitude control keeps the oscillation from growing without
bound.

\section{Active Filters and Limits}
Op-amps can implement active RC filters, but the amplifier has finite bandwidth,
slew rate, output current, noise, and stability limits. A filter that works at
audio may be inappropriate at RF.

\section{Power Supplies}
Rectifiers, filter capacitors, bleeder resistors, regulators, and switching
converters all use the same elements from the earlier chapters. Equal-value
resistors across series filter capacitors help equalize voltage and discharge
stored energy after power is removed.

\section{Exam Relevance}
\begin{exambox}
For active-circuit questions, keep three layers separate: small-signal gain,
power handling, and stability. The right answer often depends on which layer the
question is asking about.
\end{exambox}
